\begin{proof}

\begin{enumerate}[label={$\langle1\rangle$\arabic*.},leftmargin=*]
\item The local formula \CompendiumRef{II}{cdp:st:C.2.NC.1} gives
\[
 0=e(Y_t)=m\sum_i e(T_i\setminus D_i).
\]
By (C.NC.5),
\[
 e(S)-\sum_i e(T_i\setminus D_i)
 =\sum_\alpha e(C_\alpha)-2\tau=\tau.
\]
This proves \eqref{cdp:eq:C.NC.8} directly, with no base-change assumption.

\item To verify the additional comparison, locally write $t=u(z_1\cdots z_r)^m$ with $u$ nowhere zero. The normalization of $s^m=u(z_1\cdots z_r)^m$ is given by
$w^m=u$ and $s=wz_1\cdots z_r$.
The derivative of $w^m-u$ with respect to $w$ is nonzero because $w$ is a unit. Thus its map to the original chart is étale of degree $m$, and the lifted central fibre has equation $z_1\cdots z_r=0$. These local normalizations glue by the universal property of normalization; their degree and étaleness are local properties. Restriction is therefore an étale degree-$m$ covering $S'\to S$, and each triple point has exactly $m$ preimages. A lifted triangulation proves the Euler-number comparison. If the cover has several connected components, the degrees and Euler numbers add over them. The two displayed comparisons remain valid without any connectedness assertion about that cover.

\end{enumerate}
\end{proof}
